% ximeraExam: import the unmodified Q1SP17REDUX.tex activity.
% Keep this file, Q1SP17REDUX.tex, and ximeraExam.sty in the same directory.
% Run LuaLaTeX twice for each mode/output combination.
\DocumentMetadata{lang=en-US,pdfstandard=ua-2,tagging=on}
\documentclass{ximera}
\usepackage{unicode-math}
\input{preamble.tex}

% Choose points or standards mode before loading ximeraExam.
\newif\ifstandardsmode
%\standardsmodetrue
\ifstandardsmode
  \usepackage[standards]{ximeraExam}
\else
  \usepackage[points]{ximeraExam}
\fi

% Choose student or key output; do not enable both switches.
%\noprintanswers
\printanswers % Comment out when using \printanswers above.


% The imported activity uses TikZ/pgfplots; its own preamble is skipped.
% Load its graphing dependencies in the PARENT exam preamble.
\usepackage{tikz}
\usepackage{pgfplots}
\usetikzlibrary{backgrounds}



\examtitle{ximeraExam Activity Import}
\examcourse{Test Course}
\examdate{September 2026}
\examversion{Import Test}

\begin{document}
\begin{examcover}
  \examcoverheading{\examtitletext}
  \noindent Name: \rule{3in}{0.4pt}\par
  \bigskip
  \ifstandardsmode
    \standardstable[h][noearned][2]
  \else
    \gradetable[h][noearned][2]
  \fi
\end{examcover}

\begin{questions}
  % Check that numbering and metadata work before an imported activity.
  \ifstandardsmode
    \question[ALG0] Ordinary question before the import: \(1+1=\answer{2}\).
  \else
    \question[1] Ordinary question before the import: \(1+1=\answer{2}\).
  \fi

  % Import the actual activity without editing its source. Its outer exercise
  % should be ONE numbered exam question, even though it contains two nested
  % exercises, four graphical choices, and further answer/word-choice fields.
  \ifstandardsmode
    \ximeraQuestion[GRAPHS,LIMITS]{Q1SP17REDUX.tex}
  \else
    \ximeraQuestion[8]{Q1SP17REDUX.tex}
  \fi

  % Check that imported environment changes do not leak into the next question.
  \ifstandardsmode
    \question[ALG9] Ordinary question after the import: \(3+2=\answer{5}\).
  \else
    \question[2] Ordinary question after the import: \(3+2=\answer{5}\).
  \fi
\end{questions}
\end{document}
